\section{Thực nghiệm và đánh giá}

\subsection{Thiết kế kịch bản đánh giá}

Việc đánh giá hệ thống được tổ chức theo hai lớp bổ sung cho nhau:
\begin{itemize}
\item đánh giá định tính trực tiếp trên ảnh thực tế khi plugin chạy trong Paint.NET;
\item đánh giá định lượng từ notebook \texttt{dehaze-omega-regression.ipynb} trên dữ liệu cặp ảnh haze--clear.
\end{itemize}

Trong phần định tính, hai kịch bản được ưu tiên là:
\begin{itemize}
\item cảnh có haze phân bố tương đối đều trên toàn ảnh;
\item cảnh có vùng sáng mạnh hoặc vùng trời lớn, nơi fixed omega dễ xử lý quá tay.
\end{itemize}

Kịch bản thứ nhất dùng để kiểm tra khả năng phục hồi tương phản và chi tiết vùng xa của pipeline DCP cơ sở. Kịch bản thứ hai quan trọng hơn đối với adaptive omega vì đây là nhóm ảnh dễ xuất hiện hiện tượng bầu trời bị đẩy saturation hoặc mất tự nhiên nếu dùng một \(\omega\) cố định.

\subsection{Kết quả trực quan trước và sau xử lý}

Với các ảnh ngoài trời có haze mức vừa đến mạnh, ba tiêu chí quan sát quan trọng nhất là:
\begin{itemize}
\item vùng nền xa có được làm rõ hơn hay không;
\item biên vật thể có sắc nét hơn hay không;
\item màu sắc sau khôi phục có còn tự nhiên hay không.
\end{itemize}

Trong các cảnh haze phân bố khá đều, plugin cho kết quả phù hợp với mục tiêu tăng khả năng nhìn thấy của cảnh. Khi \texttt{Strength} ở mức trung bình đến khá, chi tiết vùng xa và tương phản giữa vật thể với nền thường được cải thiện rõ rệt so với ảnh gốc.

\begin{figure}[H]
  \centering
  \begin{subfigure}[t]{0.46\textwidth}
    \centering
    \includegraphics[width=0.96\textwidth]{example1.jpg}
    \caption{Ảnh hazy đầu vào}
  \end{subfigure}
  \hfill
  \begin{subfigure}[t]{0.46\textwidth}
    \centering
    \includegraphics[width=0.96\textwidth]{result_ex1.jpg}
    \caption{Kết quả sau khi khử sương mù}
  \end{subfigure}
  \vspace{0.5em}
  \caption{So sánh trước và sau xử lý trên ví dụ \texttt{example1} bằng effect DHUD DCP Dehaze}
\end{figure}

Ở ví dụ này, cải thiện dễ thấy nhất nằm ở vùng nền xa ở trung tâm lối đi và các cụm thân cây phía sau, nơi chi tiết trước đó bị phủ mờ bởi haze. Sau xử lý, tương phản giữa mặt đất, thân cây và tán lá được kéo giãn rõ hơn nhưng màu sắc tổng thể vẫn còn tương đối tự nhiên, chưa xuất hiện hiện tượng bão hòa quá mức.

Với các cảnh có vùng sáng mạnh, việc đánh giá cần thận trọng hơn. Một kết quả tốt không chỉ làm ảnh “gắt” hơn mà còn phải cân bằng được giữa việc làm rõ cảnh và tránh làm vùng trời hoặc vùng sáng trở nên giả tạo.

\begin{figure}[H]
  \centering
  \begin{subfigure}[t]{0.46\textwidth}
    \centering
    \includegraphics[width=0.96\textwidth]{example2.jpg}
    \caption{Ảnh hazy đầu vào}
  \end{subfigure}
  \hfill
  \begin{subfigure}[t]{0.46\textwidth}
    \centering
    \includegraphics[width=0.96\textwidth]{result_ex2.jpg}
    \caption{Kết quả sau khi khử sương mù}
  \end{subfigure}
  \vspace{0.5em}
  \caption{So sánh trước và sau xử lý trên ví dụ \texttt{example2} bằng effect DHUD DCP Dehaze}
\end{figure}

Trong ví dụ này, vùng mặt đường, dãy xe và khối nhà ở nửa dưới ảnh trở nên rõ biên hơn sau khi khử sương mù. Đồng thời, phần bầu trời ở nửa trên ảnh vẫn được giữ ở mức chấp nhận được, chưa bị chuyển sang tông xanh gắt quá mức. Đây là nhóm ảnh cho thấy rõ yêu cầu phải cân bằng giữa tăng độ rõ và giữ sự tự nhiên của vùng sáng rộng.

\subsection{So sánh Fixed và Adaptive}

Ở chế độ \textbf{Fixed}, toàn bộ ảnh dùng chung một \(\omega\) nên hành vi của thuật toán nhất quán nhưng kém linh hoạt. Với ảnh haze khá đồng đều, khác biệt giữa Fixed và Adaptive có thể không quá lớn nếu người dùng chọn \texttt{Strength} hợp lý.

Khác biệt trở nên rõ hơn khi ảnh chứa vùng sáng mạnh, vùng trời hoặc các mảng ít bão hòa màu. Từ công thức học được trong notebook,
\[
\omega_{\mathrm{map}}(x)=\sigma\left(-1.9900\,D(x)+1.2305\,C(x)-2.0765\,S(x)+1.8193\right),
\]
ta thấy \(\omega\) có xu hướng giảm ở nơi dark channel cao hoặc saturation thấp, và tăng ở nơi có contrast tốt hơn. Điều này giúp adaptive mode bớt mạnh tay ở những vùng dễ bị đẩy màu quá mức.

Kết quả định lượng trong notebook cũng củng cố nhận xét trên. Trên tập dữ liệu đánh giá, nhánh adaptive vượt rõ rệt baseline fixed ở cả PSNR lẫn SSIM, như thể hiện trong Bảng~\ref{tab:fixed-adaptive-metrics}.

\begin{table}[H]
  \centering
  \small
  \begin{tabular}{lcc}
    \toprule
    Chế độ & PSNR (dB) & SSIM \\
    \midrule
    Fixed (\(\omega = 0.95\)) & 17.53 & 0.8293 \\
    Adaptive omega & 24.97 & 0.9417 \\
    \bottomrule
  \end{tabular}
  \caption{So sánh định lượng giữa chế độ fixed và adaptive trên tập dữ liệu đánh giá của notebook}
  \label{tab:fixed-adaptive-metrics}
\end{table}

Các chỉ số này được lấy từ phần huấn luyện và đánh giá trên dữ liệu cặp ảnh, không phải trực tiếp từ giao diện Paint.NET. Dù vậy, chúng vẫn cho thấy việc học \(\omega(x)\) từ một số ít đặc trưng cục bộ có thể tạo ra khác biệt đáng kể so với cách dùng \(\omega\) cố định.
Kết quả này chỉ nên được dùng như bằng chứng hỗ trợ cho tính hợp lý của nhánh adaptive, chứ không đại diện tuyệt đối cho mọi ảnh thực tế mà plugin có thể gặp trong Paint.NET.

\begin{figure}[H]
  \centering
  \begin{subfigure}[t]{0.46\textwidth}
    \centering
    \includegraphics[width=0.96\textwidth]{fixed_omega.jpg}
    \caption{Kết quả với fixed omega}
  \end{subfigure}
  \hfill
  \begin{subfigure}[t]{0.46\textwidth}
    \centering
    \includegraphics[width=0.96\textwidth]{adaptive_omega.jpg}
    \caption{Kết quả với adaptive omega}
  \end{subfigure}
  \vspace{0.5em}
  \caption{So sánh fixed omega và adaptive omega trên ảnh có vùng bầu trời để thấy vai trò của adaptive trong việc tránh xử lý quá mạnh và đẩy saturation}
\end{figure}

Quan sát trực quan cho thấy ở chế độ fixed, vùng bầu trời có xu hướng bị kéo màu mạnh hơn và dễ tạo cảm giác gắt tay. Trong khi đó, adaptive giữ vùng trời dịu hơn, đồng thời vẫn cải thiện đủ độ rõ ở các cấu trúc phía dưới. Điểm mạnh của adaptive vì vậy không nằm ở việc làm ảnh “rõ hơn mọi nơi”, mà ở việc phân bổ mức dehaze hợp lý hơn giữa vùng trời và vùng cảnh.

\subsection{Đánh giá hiệu năng và trải nghiệm sử dụng}

Về trải nghiệm, điểm mạnh lớn nhất của plugin nằm ở cơ chế preview thu nhỏ kết hợp cache. Trong lúc điều chỉnh thanh \texttt{Strength}, hệ thống không phải xây lại dark channel hay transmission cho toàn ảnh lớn; preview được render từ cache trên ảnh thu nhỏ nên phản hồi đủ nhanh để hỗ trợ thao tác tương tác.

Để lượng hóa chi phí xử lý, đồ án đo thời gian trực tiếp từ lõi \texttt{DehazeProcessor} trên hai ảnh đại diện đang dùng trong báo cáo: \texttt{example1} kích thước \(600 \times 450\) và \texttt{example2} kích thước \(4400 \times 3400\). Kết quả trung bình được tách thành hai pha: xây cache (bao gồm dark channel, atmospheric light, guided filter và, nếu có, \texttt{omegaMap}) và render ảnh cuối từ cache.

\begin{table}[H]
  \centering
  \small
  \begin{tabular}{lccrrr}
    \toprule
    Ảnh & Kích thước & Chế độ & Build cache (ms) & Render (ms) & Tổng (ms) \\
    \midrule
    \texttt{example1} & \(600 \times 450\) & Fixed & 46.0 & 5.0 & 50.9 \\
    \texttt{example1} & \(600 \times 450\) & Adaptive & 55.8 & 5.0 & 60.8 \\
    \texttt{example2} & \(4400 \times 3400\) & Fixed & 2200.7 & 276.3 & 2477.0 \\
    \texttt{example2} & \(4400 \times 3400\) & Adaptive & 4075.4 & 336.1 & 4411.5 \\
    \bottomrule
  \end{tabular}
  \caption{Thời gian xử lý trung bình của lõi dehaze trên ảnh nhỏ và ảnh lớn}
  \label{tab:runtime}
\end{table}

Bảng~\ref{tab:runtime} cho thấy phần chi phí lớn nhất nằm ở bước xây cache, đặc biệt khi ảnh có độ phân giải cao hoặc khi bật adaptive omega. Ngược lại, bước render từ cache có chi phí thấp hơn đáng kể. Điều này giải thích vì sao thiết kế preview dựa trên cache là hợp lý: phần nặng được thực hiện một lần, còn thao tác thay đổi \texttt{Strength} chỉ cần đi qua pha recovery nhẹ hơn.

Mặt khác, số liệu thời gian này cũng cho thấy adaptive omega không miễn phí về mặt tính toán. Trên ảnh lớn, chi phí build cache gần như tăng mạnh do phải tính thêm đặc trưng cục bộ và suy diễn \texttt{omegaMap}. Vì vậy, adaptive phù hợp nhất như một tùy chọn tăng chất lượng cho các ảnh khó, thay vì luôn là lựa chọn tối ưu về tốc độ.

Một lợi ích thực tiễn khác là plugin hoạt động trực tiếp trong quy trình chỉnh sửa ảnh của Paint.NET. So với cách chạy thuật toán bằng script rời rồi xuất ảnh sang công cụ khác, cách tích hợp này biến dehazing thành một thao tác chỉnh sửa có thể dùng ngay trong môi trường làm việc quen thuộc.
